\begin{tabularx}{\textwidth}{@{}ll>{\raggedright\arraybackslash}X>{\raggedright\arraybackslash}X@{}}
  \toprule
  \textbf{Herramienta} & \textbf{Versión} & \textbf{Para qué se usó} & \textbf{Cómo se verificó lo obtenido} \\
  \midrule
  Claude Code & Claude Opus 5.5 & Buscar fuentes candidatas en OpenAlex, comprobar sus datos en doi.org y consultar el cuartil de cada revista en SJR. & Abrí cada DOI, descargué el artículo y comparé autores, título, revista, volumen y páginas con la ficha. Elegí las cinco fuentes y reemplacé una de las propuestas. \\
  \addlinespace
  Claude Code & Claude Opus 5.5 & Agregar mis referencias al \texttt{.bib} con el script \texttt{doi2bib.ps1}, que descarga los datos oficiales de doi.org. & El script \texttt{verificar-\allowbreak referencias.ps1} confirmó que las cinco coinciden con doi.org, y las comparé con cada artículo. \\
  \addlinespace
  Claude Code & Claude Opus 5.5 & Armar la estructura base del repositorio grupal (un archivo de fichas, uno de referencias y uno de declaración por integrante), crear las ramas de cada integrante y ajustar el formato para que las referencias muestren todos los autores y los enlaces no se salgan del margen. & Compilé el documento, revisé que no hubiera desbordes ni citas sin resolver y revisé la estructura con el grupo. \\
  \bottomrule
\end{tabularx}
